\documentclass[10pt]{article}
\usepackage{fullpage}
\usepackage{amsmath}
\usepackage{mathtools}
\usepackage{setspace}
\usepackage{booktabs}
\doublespacing

\begin{document}

\title{Cobordism Distance}
\author{Knot Theory Research}

\maketitle

\begin{abstract}

---

\end{abstract}

\pagebreak

%--INTRODUCTION--%
\section{Table}

For $K, J \subset S^3$, We define the cobordism distance as $d_C(K, J) = g_4(K\#-J)$. This collapses to a metric in $\mathcal{C}$

\begin{table}[htbp]
  \centering
  \renewcommand{\arraystretch}{1.25} % Adds breathing room between rows
  \begin{tabular}{l *{5}{c}}
    \toprule
    $d_C(K, J)$ & $O$ & $3_1$ & $4_1$ & $5_1$ & $5_2$ \\
    \midrule
    $O$         & x   & x     & x     & x     & x     \\
    $3_1$       & x   & x     & x     & x     & x     \\
    $4_1$       & x   & x     & x     & x     & x     \\
    $5_1$       & x   & x     & x     & x     & x     \\
    $5_2$       & x   & x     & x     & x     & x     \\
    \bottomrule
  \end{tabular}
  \caption{Concordance distance $d_C(K, J)$ for prime knots up to 5 crossings.}
  \label{tab:concordance-distance}
\end{table}

\subsection{Table}


%--BIBLIOGRAPHY--%
%--\begin{thebibliography}{99}--%

%--\bibitem{SampleBibLabel} Author, "Title",  Journal, Volume, Pages, Year. \\--%
%--This is the annotation for this bibliographic record.--%

%--\end{thebibliography}--%

\end{document}